%! @@TITLE@@ — Unit @@UNITINT@@, Lesson @@LESSONID@@ — Homework
% GRADUAL RELEASE: the "you do alone" block. EVERY lesson generates a homework page (user
% direction, 2026-08-31) — DeltaMath does not cover all of this course's content. Where it does,
% the teacher overrides and assigns DeltaMath instead; the plan's Reinforcement & Extension box
% names what to swap in. This page IS scored (the cover's score cell is a \blank{}, never NA).
\documentclass[10pt]{article}
\usepackage{@@PREFIX@@-article}
\usepackage{@@PREFIX@@-boxes}
% Coordinate grid: {xmin}{xmax}{ymin}{ymax}
\newcommand{\lgrid}[4]{%
  \draw[step=1, linegray!40, very thin] (#1,#3) grid (#2,#4);
  \draw[->, semithick, charcoal] (#1,0)--(#2,0) node[below right, font=\scriptsize]{$x$};
  \draw[->, semithick, charcoal] (0,#3)--(0,#4) node[left, font=\scriptsize]{$y$};}

\begin{document}

\pageheader{Unit @@UNITINT@@, Lesson @@LESSONID@@}{Homework}
% NAMESTRIP: no \namedateperiod here. The student writes their name once, on the cover.

\vspace{0.10in}

% ~6 items spanning the lesson's WHOLE standard — not a sample of it. The canonical spread:
%   1. the core procedure read off a RULE
%   2. a deliberate CONTRAST PAIR (same task, opposite condition) — the target misconception
%   3. the same procedure read off a TABLE or a GRAPH (so all three representations appear)
%   4. the SPECIAL CASE and its boundary ("for which k does this fail?")
%   5. a MODEL in a fresh context, with an interpret-the-answer follow-up
%   6. an SOL-style MULTIPLE-CHOICE item — the formative check (the plan says how to sort it)
\begin{notesbox}{Practice}
Show your work. TODO: one sentence of reminder --- the distinction or the form the lesson turned on.

\begin{enumerate}[leftmargin=1.9em, itemsep=11pt]
  \item \textbf{TODO: read the rule.} TODO \blank{1.8cm}

  \boxguard[16]
  \item \textbf{TODO: the contrast pair.} TODO
    \begin{enumerate}[label=(\alph*), leftmargin=1.9em, itemsep=4pt]
      \item TODO \blank{1.4cm}
    \end{enumerate}
    (d) TODO: the "why" item. \writelines{2}

  \item \textbf{TODO: from a table (or a graph).} TODO

  \boxguard[16]
  \item \textbf{TODO: the special case.} TODO --- and its boundary.

  \boxguard[20]
  \item \textbf{Model it.} TODO
    \begin{enumerate}[label=(\alph*), leftmargin=1.9em, itemsep=4pt]
      \item TODO: interpret the parameters. \blank{4.0cm}
      \item TODO: solve it.
        \begin{work}
          % TODO — identical in the key.
          600 + 150t &= 0 \\
                   t &= -4
        \end{work}
        TODO: does that answer mean anything here? \writelines{2}
    \end{enumerate}

  \item \textbf{Multiple choice (SOL-style).} TODO
    \begin{enumerate}[label=(\Alph*), leftmargin=2.0em, itemsep=1pt]
      \item TODO
      \item TODO
      \item TODO
      \item TODO
    \end{enumerate}
    Name one of the other three and say why it is false. \writelines{2}
\end{enumerate}
\end{notesbox}

\boxguard[18]
\begin{extensionbox}
\textbf{Build and reason.}
\begin{enumerate}[label=(\alph*), leftmargin=1.9em, itemsep=6pt]
  \item TODO: a construction --- run the lesson's procedure backwards.
        \begin{work}
          % TODO
        \end{work}
  \item TODO: a reasoning item --- "a classmate claims \dots\ explain why that cannot happen."
        \writelines{2}
\end{enumerate}
\end{extensionbox}

\begin{spiralbox}
\textbf{Coming next --- Lesson TODO: TODO Title.} TODO: two sentences connecting what they did
today to what the next lesson opens on.
\end{spiralbox}

\end{document}
